% TOP_PROBLEM: {"id": 30000755, "title": "Single-Solution Bounds for Diophantine Equations", "category_id": 1, "category": {"id": 1, "name": "number_theory", "display_name": "Number Theory", "description": "Properties of integers, prime numbers, Diophantine equations.", "slug": "number-theory", "order_index": 1, "created_at": "2026-07-31T15:26:25.670Z"}, "status": "open"}
% FIRST_POSTED: 2026-10-06
% ATTEMPT_STATUS: unresolved
% OPENAI_CLAIM: {"schema_version":1,"source_repo":"https://github.com/openai/math","source_commit":"adc7f1241b42e322a6451854ab7e4b4c146bf78a","release_date":"2026-10-06","match":"partial","resolution":"partial","independently_reviewed":false,"second_pass":"n/a","families":[{"family":"242","title":"Single-fold Diophantine representations and undecidability under an at-most-one-solution promise","subject":"Mathematical logic","manuscripts":[{"title":"Single-fold Diophantine representations","pdf_path":"preprints/Single-fold-Diophantine-representations-September-24-2026/paper.pdf","readme_path":"preprints/Single-fold-Diophantine-representations-September-24-2026/README.md","date":"2026-09-24","theorem_ref":"thm:main (main.tex:56)","statement_sources":[{"path":"preprints/Single-fold-Diophantine-representations-September-24-2026/build/main.tex","line":56,"label":"thm:main"}]}],"lean":{"doc_path":"lean/docs/242.md","comparators":["lean/ComparatorChallenges/SingleFold.json","lean/ComparatorChallenges/SingleFold.lean"],"declarations":["OAI.SingleFold.main"],"covers_main_theorem":true},"reasoning_trace":null}]}
\documentclass[11pt]{article}
\usepackage[margin=25mm]{geometry}
\usepackage{fontspec}
\setmainfont{Latin Modern Roman}
\usepackage{amsmath,amssymb,mathtools}
\usepackage{unicode-math}
\setmathfont{Latin Modern Math}
\usepackage{enumitem}
\usepackage{xurl,hyperref,bookmark}
\hypersetup{colorlinks=true,urlcolor=blue}
\setcounter{secnumdepth}{0}
\setlength{\parindent}{0pt}
\setlength{\parskip}{5pt}
\setlength{\emergencystretch}{3em}
\begin{document}
\section*{Single-Solution Bounds for Diophantine Equations}

\subsection{OpenAI's claim}
OpenAI presents the following claims relevant to this question, with manuscript numbers following paper order in Sources. In manuscript 1 (family 242; 2026-09-24), the asserted result is that every recursively enumerable relation on natural-number tuples has an integer-polynomial representation with exactly one complete auxiliary witness for each member.

\subsection{Mathematical statement of the claimed result}
\paragraph{Manuscript 1.}
Use $\mathbb N$ including zero. For each $n\ge1$ and recursively enumerable $S\subseteq\mathbb N^n$, the claim gives an integer $m\ge1$ and an integral polynomial $P(a,w)$ with
\[
 \forall a\in\mathbb N^n:\quad
 |\{w\in\mathbb N^m:P(a,w)=0\}|=
 \begin{cases}1,&a\in S,\\0,&a\notin S.\end{cases}
\]
The number of witness variables is fixed for the set $S$. This asserts a unique Diophantine witness for every member, rather than merely a finite number of witnesses.

\subsection{Outline of the argument}
For manuscript 1 (family 242; 2026-09-24), pell equations recover exponentiation, but their indices are selected using a uniquely coded rank-one elliptic curve. Bounded height error gives an exact growth cutoff, and canonical arithmetic codes then encode each halting computation with a unique full tuple.

\subsection{Scope relative to this problem}
The theorem has unique natural-number witnesses, so undecidable recursively enumerable sets preclude a computable bound for all such witnesses. The problem's witness domain is not specified precisely enough to identify every requested single-solution clause with that theorem, and complete coverage is not established here.

\subsection{Formal verification}
For manuscript 1 (family 242; 2026-09-24), OpenAI lists Lean statements for this assertion: every recursively enumerable relation on natural-number tuples has an integer-polynomial representation with exactly one complete auxiliary witness for each member.

The selected statement includes the clause used for this problem.

The linked family Comparator checks are \texttt{SingleFold.lean}.

Their named declarations include \texttt{OAI.SingleFold.main}.

This site has not executed the Lean code. Described statement coverage is not an independent proof check.

\subsection{Status caveat}
The manuscripts were produced by an internal OpenAI model. OpenAI cautions that results without formalization can contain mistakes. This site has not independently checked the proofs; the label records the scope of OpenAI's claim.

\subsection{Sources}
\begin{itemize}
\item OpenAI, \emph{Single-fold Diophantine representations}, 2026-09-24; citation key \texttt{OAI:Single-fold-Diophantine-representations-September-24-2026}. \href{https://github.com/openai/math/blob/adc7f1241b42e322a6451854ab7e4b4c146bf78a/preprints/Single-fold-Diophantine-representations-September-24-2026/paper.pdf}{PDF (GitHub)}; \href{https://github.com/openai/math/raw/adc7f1241b42e322a6451854ab7e4b4c146bf78a/preprints/Single-fold-Diophantine-representations-September-24-2026/paper.pdf}{PDF (direct)}; \href{https://github.com/openai/math/blob/adc7f1241b42e322a6451854ab7e4b4c146bf78a/preprints/Single-fold-Diophantine-representations-September-24-2026/README.md}{Manuscript page and citation}; \href{https://github.com/openai/math/blob/main/preprints/Single-fold-Diophantine-representations-September-24-2026/paper.pdf}{Latest version}; Statement source: \href{https://github.com/openai/math/blob/adc7f1241b42e322a6451854ab7e4b4c146bf78a/preprints/Single-fold-Diophantine-representations-September-24-2026/build/main.tex\#L56}{thm:main}.
\item Family 242: \href{https://github.com/openai/math/blob/adc7f1241b42e322a6451854ab7e4b4c146bf78a/lean/docs/242.md}{Lean scope}.
\item \href{https://github.com/openai/math/blob/adc7f1241b42e322a6451854ab7e4b4c146bf78a/lean/ComparatorChallenges/SingleFold.json}{Comparator SingleFold.json}.
\item \href{https://github.com/openai/math/blob/adc7f1241b42e322a6451854ab7e4b4c146bf78a/lean/ComparatorChallenges/SingleFold.lean}{Comparator SingleFold.lean}.
\item \href{https://github.com/openai/math/blob/adc7f1241b42e322a6451854ab7e4b4c146bf78a/CONTENTS.md}{OpenAI catalogue entry}.
\end{itemize}
\end{document}
